% =====================================================================
%  MODELO EM BRANCO -- uma atividade do MOP em formato metodológico
%
%  Como usar:
%   1. Copie este arquivo para capitulos/<subcomponente>/04-x-<nome>.tex
%   2. Substitua tudo o que estiver entre << >>
%   3. Inclua o arquivo no capitulo.tex do subcomponente com \input{...}
%
%  Regras de redação (ver "Apresentação do modelo"):
%   - texto impessoal, no futuro do presente ("serão realizadas");
%   - informação textual em quadro; dados numéricos em tabela;
%   - toda figura, quadro e tabela com \fonte{...};
%   - inconsistências do MOP registradas com \pendencia{...}.
% =====================================================================
\subsection{Atividade <<X.X.X.X>> -- <<Nome da atividade, com verbo no infinitivo>>}
\label{subsec:<<rotulo-curto>>}

% ---------------------------------------------------------------------
\subsubsection{Objetivo e fundamentação}
% O que a atividade entrega (1 frase) + por que é necessária (1 a 2
% parágrafos) + premissas ou dependências de partida.
% ---------------------------------------------------------------------

<<Objetivo da atividade.>>

<<Fundamentação: problema que a atividade resolve e referências
normativas ou técnicas que a sustentam \cite{mop2026}.>>

% ---------------------------------------------------------------------
% Ficha-síntese: manter TODOS os campos e nesta ordem.
% Quando o MOP não trouxer a informação, escrever "A definir" e
% registrar uma \pendencia{...}.
% ---------------------------------------------------------------------
\begin{fichaatividade}{Ficha-síntese da Atividade <<X.X.X.X>>}{qua:ficha-<<rotulo-curto>>}
  \campo{Código}{<<X.X.X.X>>}
  \campo{Tipo de atividade}{<<Estudo e mobilização | Serviço e obra | Aquisição | ...>>}
  \campo{Forma de execução}{<<Contratação pública (licitação) | Convênio | Execução direta>>}
  \campo{Fonte de recursos}{<<BIRD | Contrapartida | Tesouro estadual>>}
  \campo{Responsável}{<<Instituição executora>>}
  \campo{Dependências institucionais}{<<Instituições das quais a execução depende>>}
  \campo{Prazo estimado}{<<Ano X ao ano Y>>}
  \campo{Produto esperado}{<<Produto final, com quantidade>>}
  \campo{Unidade de medida}{<<Unidade>>}
  \campo{Evidência de comprovação}{<<Documento que comprova a entrega>>}
  \campo{Periodicidade}{<<Única | Por lote | Semestral | Contínua>>}
  \campo{Validação}{<<Instituição que valida>>}
\end{fichaatividade}

% ---------------------------------------------------------------------
\subsubsection{Procedimentos}
% Agrupar os passos em FASES (\paragraph). Em cada passo: nome em
% negrito + o que será feito. Usar "resume" para continuar a numeração
% das alíneas entre fases.
% ---------------------------------------------------------------------

\paragraph{Fase 1 -- <<Nome da fase>>.}
\begin{enumerate}[label=\alph*)]
  \item \textbf{<<passo>>}: <<descrição do que será feito>>;
  \item \textbf{<<passo>>}: <<descrição>>.
\end{enumerate}

\paragraph{Fase 2 -- <<Nome da fase>>.}
\begin{enumerate}[label=\alph*), resume]
  \item \textbf{<<passo>>}: <<descrição>>.
\end{enumerate}

% Opcional: fluxograma da atividade (ver figuras/fluxograma-esgotamento.tex)
% \begin{figure}[htbp]
%   \centering
%   \caption{Fluxograma da Atividade <<X.X.X.X>>}
%   \label{fig:fluxograma-<<rotulo-curto>>}
%   \input{figuras/<<arquivo>>}
%   \fonte{Elaborado pelo IAT (2026).}
% \end{figure}

% ---------------------------------------------------------------------
\subsubsection{Etapas, produtos e evidências}
% Quadro detalhado por etapa, em página paisagem. Mesmas colunas em
% todas as atividades.
% ---------------------------------------------------------------------

\begin{landscape}
\begin{quadro}[p]
\caption{Etapas da Atividade <<X.X.X.X>>}
\label{qua:etapas-<<rotulo-curto>>}
\scriptsize
\renewcommand{\arraystretch}{1.25}
\begin{tabularx}{\linewidth}{|P{2.6cm}|L|P{2.1cm}|Q{1.3cm}|P{2.5cm}|P{2cm}|P{2.6cm}|Q{1.4cm}|P{1.7cm}|}
\hline
\cab{Etapa} & \cab{Descrição técnica} & \cab{Responsável} & \cab{Prazo} &
\cab{Produto esperado} & \cab{Unidade} & \cab{Evidência} & \cab{Periodic.} &
\cab{Validação} \\ \hline
1. <<Etapa>> & 1.1 <<subetapa>>; 1.2 <<subetapa>> & <<>> & <<Ano X a Y>> &
<<>> & <<>> & <<>> & <<>> & <<>> \\ \hline
2. <<Etapa>> & 2.1 <<subetapa>> & <<>> & <<>> &
<<>> & <<>> & <<>> & <<>> & <<>> \\ \hline
\end{tabularx}
\fonte{Adaptado de \citeonline{mop2026}.}
\end{quadro}
\end{landscape}

% ---------------------------------------------------------------------
\subsubsection{Interfaces}
% Com quais atividades, subcomponentes ou indicadores esta atividade se
% relaciona, e em que sentido (recebe insumo de / entrega produto para).
% ---------------------------------------------------------------------

\begin{itemize}
  \item \atividade{<<X.X.X.X>>}: <<relação>>.
  \item Indicador <<nome do indicador>>: <<contribuição>>.
\end{itemize}
